\section{Discussão}

\subsection{Interpretação dos Resultados}

Para RQ1, GraphQL apresentou menor tempo de resposta médio (462,76 ms)
em comparação com REST (651,58 ms). O teste de Wilcoxon-Mann-Whitney
rejeitou $H_0$ ($p < 0,001$), indicando que GraphQL é significativamente
mais rápido que REST para a query analisada. Uma possível explicação é que
a API REST do GitHub retorna o objeto completo do repositório com dezenas
de campos independentemente do repositório consultado, enquanto a query
GraphQL seleciona apenas três campos, reduzindo o tempo de serialização
e transferência.

Para RQ2, GraphQL apresentou payloads significativamente menores
($p < 0,001$), com mediana de 82 bytes (variação de 81 a 87)
contra 6277 bytes do REST. Isso confirma a vantagem fundamental do
GraphQL: o cliente especifica exatamente os campos desejados.

\subsection{Ameaças à Validade}

O experimento utilizou 5 repositórios do GitHub como objetos experimentais,
o que reduz o viés de seleção em comparação com um único repositório,
porém ainda está limitado a uma única plataforma de API (validade externa).
APIs de outros provedores (Google, Twitter, etc.) podem apresentar
comportamentos diferentes.

A query GraphQL utilizada é intencionalmente simples (3 campos por recurso);
consultas mais complexas, com \emph{nested queries} ou múltiplos recursos,
podem apresentar comportamento diferente em ambas as interfaces.

Recomenda-se a replicação do experimento com:
\begin{itemize}
    \item APIs não-GitHub (ex.: SWAPI, Spotify, Cloudflare)
    \item Consultas com maior profundidade e cardinalidade
    \item Medição de impacto no servidor (CPU e memória)
\end{itemize}

\subsection{Trabalhos Relacionados}

Os resultados estão alinhados com \cite{brito2019graphql}, que encontraram
vantagens do GraphQL no tamanho das respostas, mas alertaram para possíveis
impactos negativos no tempo de resposta em consultas complexas. O artigo
original do GraphQL \cite{lee2015graphql} já apontava a seletividade de campos
como benefício central, confirmado pelos nossos dados. Maiores detalhes sobre
os endpoints utilizados podem ser encontrados na documentação oficial
\cite{github2024api}.
